% Winkel und Dreiecke – bank/winkel-dreiecke/ (zusammenbau v0.4, Heft)
\einheitenkopf[t7]{Winkel und Dreiecke}
\setcounter{aufgabe}{7}

\begin{aufgabe}{Winkel an Geraden – Prüfungsaufgaben}
\begin{teile}
\teil Die Geraden g und h schneiden sich; zwischen ihnen liegt $\alpha = 64^\circ$. Eine dritte Gerade durch den Schnittpunkt teilt den Scheitelwinkel von α in zwei Teile, der untere misst $23^\circ$. Wie groß ist der obere Teil β? \hfill (P10 2014 OS) \\ β = \leerfeld °
\teil $g \parallel h$, g liegt oben, h unten, die Gerade k schneidet beide. An g liegt unterhalb von g links von k ein Winkel von $49^\circ$. Wie groß ist der Winkel an h oberhalb von h rechts von k? \hfill (P10 2020 OS) \\

\parallelenpaar{49}{}{}{}{}
\leerfeld °
\teil Im Dreieck ABC ist der Winkel bei C $58^\circ$ groß. Die Seiten AB und CB werden über B hinaus verlängert; zwischen den beiden Verlängerungen liegt ein Winkel von $43^\circ$. Wie groß ist der Innenwinkel β bei B? \hfill (P10 2019 OS) \\ β = \leerfeld °
\teil Im Trapez ABCD mit $AB \parallel CD$ ist $\beta = 58^\circ$, $\gamma = 122^\circ$ und $\delta = 107^\circ$. Wie groß ist α? \hfill (P10 2021 OS) \\

\viereck[winkel]{(0,0)}{(6,0)}{(4.44,2.5)}{(0.76,2.5)}
α = \leerfeld °
\end{teile}
\end{aufgabe}

\begin{aufgabe}{Winkel an Geraden – Prüfungsaufgaben – weiter}
\begin{teile}
\teil Die parallelen Geraden h (oben) und g (unten) werden von der Geraden k geschnitten. An h liegt unterhalb von h links von k ein Winkel von $44^\circ$; α liegt an g oberhalb von g links von k. Wie groß ist α? \hfill (P10 2015 OS) \\

\parallelenpaar{44}{}{}{}{}
α = \leerfeld °
\teil Eine gerade Straße k überquert zwei parallele Gleise, in der Skizze liegt Gleis h oben, Gleis g unten. Unterhalb von h links von k bilden Straße und Gleis einen Winkel von $51^\circ$. Wie groß ist der Winkel α an g oberhalb von g links von k? \hfill (P10 2015 OS) \\

\parallelenpaar{51}{}{}{}{}
α = \leerfeld °
\teil Ein Parallelogramm hat links unten einen Winkel von $79^\circ$. Wie groß ist der benachbarte Winkel β rechts unten? \hfill (P10 2026 FOR) \\

\parallelogramm[winkel={79^\circ,\beta,,}]{4}{2.5}{79}
β = \leerfeld °
\teil Ein Scherengitter ist aus Parallelogrammen aufgebaut. In einem davon misst der Winkel links unten $83^\circ$. Wie groß ist der benachbarte Winkel β rechts unten? \hfill (P10 2026 FOR) \\

\parallelogramm[winkel={83^\circ,\beta,,}]{4}{2.5}{83}
β = \leerfeld °
\end{teile}
\end{aufgabe}

\begin{aufgabe}{Winkelsumme – Prüfungsaufgaben}
\begin{teile}
\teil Ergänze richtig: „In jedem Drachenviereck …“ \hfill (P10 2026 FOR) \\ \kreuz{sind alle Winkel gleich groß} \\ \kreuz{gibt es zwei Paare gleich langer Nachbarseiten} \\ \kreuz{sind alle Seiten gleich lang}
\teil Ergänze richtig: „In jeder Raute …“ \hfill (P10 2016 OS) \\ \kreuz{sind benachbarte Winkel gleich groß} \\ \kreuz{sind gegenüberliegende Winkel gleich groß} \\ \kreuz{sind alle Winkel gleich groß}
\teil Im Dreieck ABC ist $\alpha = \beta = 66^\circ$ und AC = 7 cm. Gib die Länge von BC und die Größe von γ an. \hfill (P10 2023 OS) \\ BC = \leerfeld[cm] , γ = \leerfeld °
\teil Die Giebelwand eines Zeltes ist ein Dreieck ABC mit den Winkeln $\alpha = \beta = 73^\circ$ an den Bodenecken; die Kante AC ist 2,4 m lang. Wie lang ist die Kante BC, und wie groß ist der Winkel γ an der Spitze? \hfill (P10 2023 OS) \\ BC = \leerfeld[m] , γ = \leerfeld °
\end{teile}
\end{aufgabe}

\begin{aufgabe}{Konstruieren – Prüfungsaufgaben}
\begin{teile}
\teil Ein dreieckiges Beet ABC hat die Seite a = 12 m. Was gilt für die Summe der beiden anderen Seiten b und c? \hfill (P10 2020 OS) \\ \kreuz{$b + c > 12$ m} \\ \kreuz{$b + c = 12$ m} \\ \kreuz{$b + c < 12$ m}
\teil Drei Holzstäbe sollen zu einem Dreieck gelegt werden; der längste Stab a ist 15 cm lang. Was muss für die beiden anderen Stäbe b und c gelten? \hfill (P10 2020 OS) \\ \kreuz{$b + c > 15$ cm} \\ \kreuz{$b + c = 15$ cm} \\ \kreuz{$b + c < 15$ cm}
\end{teile}
\end{aufgabe}
